\documentclass[11pt,a4paper]{article}

% ===================== POLICES & LANGUE =====================
\usepackage{fontspec}
\setmainfont{Latin Modern Roman}
\setsansfont{Latin Modern Sans}
\setmonofont{Latin Modern Mono}
\usepackage{polyglossia}
\setdefaultlanguage{french}

% ===================== MISE EN PAGE =====================
\usepackage[margin=2.5cm]{geometry}
\usepackage{amsmath,amssymb,graphicx,xcolor,booktabs,hyperref,tcolorbox,enumitem,fancyhdr,array,pifont,longtable,multirow,listings}
\tcbuselibrary{theorems,breakable,skins}
\usepackage{tikz}
\usetikzlibrary{shapes.geometric,arrows.meta,positioning,calc,backgrounds,fit}

\newcommand{\cmark}{\ding{51}}
\newcommand{\xmark}{\ding{55}}
\hypersetup{colorlinks=true,linkcolor=blue!60!black,urlcolor=blue!60!black,citecolor=blue!60!black}

% ===================== BOÎTES =====================
\newtcolorbox{objectifbox}[1][]{colback=blue!5,colframe=blue!50,fonttitle=\bfseries,title={#1},breakable}
\newtcolorbox{infobox}[1][]{colback=green!5,colframe=green!50!black,fonttitle=\bfseries,title={#1},breakable}
\newtcolorbox{attentionbox}[1][]{colback=orange!5,colframe=orange!70,fonttitle=\bfseries,title={#1},breakable}
\newtcolorbox{retenirbox}[1][]{colback=yellow!10,colframe=yellow!50!black,fonttitle=\bfseries,title={#1},breakable}
\newtcolorbox{prereqbox}[1][]{colback=cyan!5,colframe=cyan!50!black,fonttitle=\bfseries,title={#1},breakable}
\newtcolorbox{projetbox}[1][]{colback=teal!5,colframe=teal!50!black,fonttitle=\bfseries,title={#1},breakable}

% ===================== CODE (listings) =====================
\definecolor{codebg}{gray}{0.96}
\definecolor{kwcolor}{RGB}{0,90,160}
\definecolor{strcolor}{RGB}{160,55,30}
\definecolor{comcolor}{RGB}{95,120,95}
\definecolor{ndcolor}{RGB}{120,80,150}

\lstdefinelanguage{JS}{
  morekeywords={let,const,var,function,return,if,else,for,of,in,while,do,true,false,null,undefined,typeof,break,continue,new,this},
  morekeywords=[2]{console,log,table,document,Math,window,alert,prompt,forEach,map,filter,push,pop,querySelector,querySelectorAll,addEventListener,classList,textContent,innerHTML,createElement,appendChild,getElementById},
  sensitive=true,
  morecomment=[l]{//},
  morecomment=[s]{/*}{*/},
  morestring=[b]",
  morestring=[b]',
  morestring=[b]`,
}

\lstdefinestyle{js}{
  language=JS,
  basicstyle=\ttfamily\small,
  keywordstyle=\color{kwcolor},
  keywordstyle=[2]\color{ndcolor},
  commentstyle=\color{comcolor}\itshape,
  stringstyle=\color{strcolor},
  backgroundcolor=\color{codebg},
  frame=single,
  rulecolor=\color{gray!45},
  breaklines=true,
  breakatwhitespace=false,
  showstringspaces=false,
  numbers=none,
  tabsize=2,
  columns=fullflexible,
  keepspaces=true,
  literate=%
    {é}{{é}}1 {è}{{è}}1 {à}{{à}}1 {â}{{â}}1 {ê}{{ê}}1
    {î}{{î}}1 {ô}{{ô}}1 {û}{{û}}1 {ù}{{ù}}1 {ç}{{ç}}1
    {ï}{{ï}}1 {ë}{{ë}}1 {É}{{É}}1 {È}{{È}}1 {À}{{À}}1
    {œ}{{œ}}1 {°}{{°}}1 {²}{{²}}1
}

\lstdefinestyle{html}{
  basicstyle=\ttfamily\small,
  backgroundcolor=\color{codebg},
  frame=single,
  rulecolor=\color{gray!45},
  breaklines=true,
  showstringspaces=false,
  numbers=none,
  tabsize=2,
  columns=fullflexible,
  keepspaces=true,
  literate=%
    {é}{{é}}1 {è}{{è}}1 {à}{{à}}1 {ê}{{ê}}1 {ç}{{ç}}1 {î}{{î}}1
}

\lstset{style=js}
\newcommand{\code}[1]{\texttt{#1}}

% ===================== EN-TÊTES =====================
\pagestyle{fancy}
\fancyhf{}
\fancyhead[L]{\small Jàng Afrig --- Lab 5 : JavaScript en pratique}
\fancyhead[R]{\small UCAD --- L1 DUT1-Info}
\fancyfoot[C]{\thepage}
\fancyfoot[L]{\small Dr. El Hadji Bassirou TOURÉ}
\fancyfoot[R]{\small 2025--2026}
\renewcommand{\headrulewidth}{0.4pt}
\renewcommand{\footrulewidth}{0.4pt}

% ===================== TITRE =====================
\title{%
\vspace{-1cm}
\textbf{\LARGE Fondamentaux du Web}\\
\textbf{\LARGE HTML, CSS, JavaScript}\\[0.8cm]
\textbf{\huge Lab 5 --- JavaScript en pratique}\\[0.3cm]
\textbf{\Large Variables, fonctions, conditions}\\[0.5cm]
\large Licence 1 DUT1-Informatique --- 2025--2026}
\date{}
\author{}

\begin{document}
\maketitle
\thispagestyle{empty}

\begin{prereqbox}[Informations pratiques]
\textbf{Où sommes-nous ?} Au CM3, on a découvert le langage JavaScript dans la \emph{console} (notre brouillon). Aujourd'hui, on quitte le brouillon : on écrit notre \textbf{premier vrai fichier de code}, \code{script.js}, branché à la page Jàng Afrig.

\smallskip
\textbf{Avant de commencer, vous devez avoir :}
\begin{itemize}[leftmargin=*,noitemsep]
  \item Votre repo \code{jang-afrig} ouvert dans VS Code, avec Live Server qui fonctionne.
  \item Les pages HTML et le CSS des Labs précédents (jusqu'au tag \code{v0.4}).
  \item Les notions du CM3 en tête : \code{let}/\code{const}, types, conditions, boucles, fonctions, et les \textbf{trois règles d'or}.
\end{itemize}

\smallskip
\textbf{Validation de fin de séance :} votre fichier \code{script.js} contient toutes les fonctions d'entraînement, la console les exécute \emph{sans erreur rouge}, et vous avez poussé un commit avec le tag \code{v0.5} ainsi qu'un rapport.
\end{prereqbox}

\begin{objectifbox}[Objectifs du Lab 5]
À la fin de cette séance, vous saurez :
\begin{enumerate}[leftmargin=*,noitemsep]
  \item Créer un fichier \code{script.js} et le brancher à \code{index.html}.
  \item Afficher des messages dans la console avec \code{console.log}.
  \item Déclarer des variables et écrire des fonctions qui renvoient un résultat.
  \item Prendre des décisions avec \code{if}/\code{else} et le ternaire.
  \item Parcourir un tableau avec \code{for...of}.
  \item Coder cinq fonctions utiles : convertisseur FCFA, moyenne, majorité, FizzBuzz « sénégalais », palindrome.
\end{enumerate}
\end{objectifbox}

\tableofcontents
\bigskip\hrule\bigskip

\section*{Rappel express du CM3}
\addcontentsline{toc}{section}{Rappel express du CM3}
Trois choses à garder sous la main :
\begin{itemize}[leftmargin=*,noitemsep]
  \item \textbf{Variables} : \code{const} par défaut, \code{let} si ça change, jamais \code{var}.
  \item \textbf{Comparaison} : toujours \code{===}, jamais \code{==}.
  \item \textbf{Fonction} : \code{function nom(parametres) \{ ... return ...\, \}}.
\end{itemize}

% ===================== SECTION 1 =====================
\section{Mise en place : créer et brancher \texttt{script.js}}

\subsection{Créer le fichier}
Dans VS Code, à la racine de votre dossier \code{jang-afrig} (au même endroit que \code{index.html} et \code{style.css}), créez un nouveau fichier nommé exactement \code{script.js}.

\subsection{Le brancher à la page}
Pour que la page exécute votre code, on ajoute la balise \code{<script>} dans le \code{<head>} de \code{index.html}, avec le mot \code{defer} (rappel du CM3 : « attends que la page soit prête »).

\begin{lstlisting}[style=html]
<head>
  <meta charset="UTF-8">
  <title>Jàng Afrig</title>
  <link rel="stylesheet" href="style.css">
  <script src="script.js" defer></script>
</head>
\end{lstlisting}

\subsection{Vérifier que tout est branché}
Dans \code{script.js}, écrivez une seule ligne :
\begin{lstlisting}[style=js]
console.log("script.js est bien branché !");
\end{lstlisting}

\begin{infobox}[Vérification]
Lancez la page avec Live Server, puis appuyez sur \textbf{F12} et ouvrez l'onglet \textbf{Console}. Vous devez voir le message \emph{« script.js est bien branché ! »}. Si rien n'apparaît, vérifiez l'orthographe du nom de fichier et la balise \code{<script>}. \textbf{Si ça ne marche pas, levez la main.}
\end{infobox}

\begin{attentionbox}[La page ne « change » pas encore, c'est normal]
Aujourd'hui, tout votre code s'affiche dans la \textbf{console}, pas sur la page. C'est voulu : apprendre à modifier la page visible viendra au \textbf{Lab 7} (le DOM). Pour l'instant, la console est votre terrain d'entraînement.
\end{attentionbox}

% ===================== SECTION 2 =====================
\section{Parler à la console}

On commence doux. Tapez ceci dans \code{script.js} (remplacez par votre prénom) :

\begin{lstlisting}[style=js]
const prenom = "Fatou";
console.log("Bonjour " + prenom);
console.log(`Bonjour ${prenom}, prêt(e) à coder ?`);
\end{lstlisting}

Les deux dernières lignes affichent un message personnalisé. La première version \emph{colle} le texte avec \code{+} ; la seconde utilise une \emph{template string} (les accents graves) du CM3 : plus lisible.

\begin{infobox}[Vérification]
La console doit afficher deux lignes commençant par \emph{« Bonjour »}. Changez la valeur de \code{prenom}, sauvegardez : Live Server recharge, le message change.
\end{infobox}

% ===================== SECTION 3 =====================
\section{Variables et types, pour de vrai}

Rangeons quelques informations sur un joueur de Jàng Afrig :

\begin{lstlisting}[style=js]
const pseudo = "Aminata";   // texte (string)
let score = 0;              // nombre (number), va changer -> let
const partieFinie = false;  // booléen (boolean)

console.log(`Joueur : ${pseudo}`);
console.log(`Score de départ : ${score}`);

score = score + 5;          // on réécrit l'ardoise
console.log(`Nouveau score : ${score}`);  // 5
\end{lstlisting}

\begin{infobox}[Vérification]
La console doit montrer le score passer de \code{0} à \code{5}. Essayez ensuite de réaffecter \code{pseudo} (\code{pseudo = "Moussa";}) : la console affichera une \textbf{erreur rouge}, car \code{pseudo} est un \code{const}. C'est exactement ce qu'on attend : c'est la preuve que l'étiquette est « gravée ».
\end{infobox}

% ===================== SECTION 4 =====================
\section{Ma première fonction}

Une fonction est une machine réutilisable (CM3). Écrivons-en une qui fabrique un message de bienvenue :

\begin{lstlisting}[style=js]
function direBonjour(nom) {
  return `Bonjour ${nom}, bienvenue sur Jàng Afrig !`;
}

console.log(direBonjour("Cheikh"));
console.log(direBonjour("Ndèye"));
\end{lstlisting}

On a écrit la machine \emph{une fois}, et on l'appelle \emph{autant de fois} qu'on veut, avec des noms différents.

\begin{objectifbox}[Exercice éclair]
Écrivez une fonction \code{direAuRevoir(nom)} qui renvoie \emph{« À bientôt, [nom] ! »}, et appelez-la deux fois avec deux prénoms.
\end{objectifbox}

% ===================== SECTION 5 =====================
\section{Prendre des décisions : les conditions}

Une fonction peut \textbf{choisir} quoi renvoyer selon la situation. Reprenons l'appréciation d'une note (CM3) :

\begin{lstlisting}[style=js]
function appreciation(note) {
  if (note >= 16) {
    return "Très bien";
  } else if (note >= 12) {
    return "Bien";
  } else if (note >= 10) {
    return "Passable";
  } else {
    return "À retravailler";
  }
}

console.log(appreciation(14)); // "Bien"
console.log(appreciation(8));  // "À retravailler"
\end{lstlisting}

Quand le choix est binaire, le \textbf{ternaire} fait la même chose en une ligne :

\begin{lstlisting}[style=js]
function statutMajorite(age) {
  return age >= 18 ? "majeur" : "mineur";
}

console.log(statutMajorite(20)); // "majeur"
\end{lstlisting}

% ===================== SECTION 6 =====================
\section{Répéter sans se fatiguer : \texttt{for...of}}

Pour traiter \emph{chaque} élément d'une liste, on utilise \code{for...of} (CM3). Affichons une petite liste de pays :

\begin{lstlisting}[style=js]
const pays = ["Sénégal", "Mali", "Ghana", "Kenya"];

for (const nom of pays) {
  console.log(`Pays africain : ${nom}`);
}

console.log(`Il y a ${pays.length} pays dans la liste.`);
\end{lstlisting}

Notez \code{pays.length} : la \emph{longueur} du tableau (ici 4). On s'en servira beaucoup.

\begin{infobox}[Vérification]
La console doit afficher quatre lignes « Pays africain : \ldots », puis « Il y a 4 pays dans la liste. ».
\end{infobox}

% ===================== EXERCICES =====================
\section{Exercices d'application}

C'est le cœur de la séance. Ajoutez ces fonctions à \code{script.js} et testez-les avec des \code{console.log}.

\begin{objectifbox}[Exercice 1 --- Convertisseur FCFA $\rightarrow$ EUR]
Le taux est fixe : \textbf{1 EUR = 656 FCFA}. Écrivez \code{fcfaVersEuro(montant)} qui renvoie le montant en euros.\\
\emph{Indice :} il faut \textbf{diviser} par 656. Testez avec \code{10000} FCFA.
\end{objectifbox}

\begin{lstlisting}[style=js]
function fcfaVersEuro(montant) {
  return montant / 656;
}

console.log(fcfaVersEuro(10000)); // environ 15.24
\end{lstlisting}

\begin{objectifbox}[Exercice 2 --- Moyenne d'un étudiant]
Écrivez \code{moyenne(notes)} qui reçoit un tableau de notes et renvoie leur moyenne.\\
\emph{Indice :} additionnez toutes les notes dans une variable \code{somme} avec \code{for...of}, puis divisez par \code{notes.length}.
\end{objectifbox}

\begin{lstlisting}[style=js]
function moyenne(notes) {
  let somme = 0;
  for (const note of notes) {
    somme = somme + note;
  }
  return somme / notes.length;
}

const notesAminata = [12, 15, 9, 14];
console.log(`Moyenne : ${moyenne(notesAminata)}`); // 12.5
\end{lstlisting}

\begin{objectifbox}[Exercice 3 --- Majorité]
Écrivez \code{estMajeur(age)} qui renvoie un \textbf{booléen} (\code{true} ou \code{false}).\\
\emph{Indice :} la comparaison \code{age >= 18} est \emph{déjà} un booléen ! On peut la renvoyer directement.
\end{objectifbox}

\begin{lstlisting}[style=js]
function estMajeur(age) {
  return age >= 18;
}

console.log(estMajeur(20)); // true
console.log(estMajeur(16)); // false
\end{lstlisting}

\subsection*{Exercice 4 --- FizzBuzz « à la sénégalaise »}
\addcontentsline{toc}{subsection}{Exercice 4 --- FizzBuzz sénégalais}

Un classique de la programmation, version Jàng Afrig. Pour chaque nombre de 1 à 30 :
\begin{itemize}[leftmargin=*,noitemsep]
  \item s'il est multiple de \textbf{3 et de 5} (donc de 15) $\rightarrow$ afficher \emph{« Thiéboudienne »} ;
  \item sinon, s'il est multiple de \textbf{3} $\rightarrow$ afficher \emph{« Thié »} ;
  \item sinon, s'il est multiple de \textbf{5} $\rightarrow$ afficher \emph{« Bou »} ;
  \item sinon $\rightarrow$ afficher le nombre.
\end{itemize}

\textbf{Le piège : l'ordre des tests compte.} Si on testait « multiple de 3 » \emph{avant} « multiple de 15 », le nombre 15 afficherait « Thié » au lieu de « Thiéboudienne ». On teste donc \textbf{toujours le cas le plus précis en premier} (le 15).

\begin{center}
\begin{tikzpicture}[
  font=\footnotesize,
  deci/.style={diamond, aspect=2.4, draw, thick, align=center, inner sep=1pt, fill=cyan!6, draw=cyan!55!black},
  io/.style={rectangle, rounded corners=3pt, draw, thick, align=center, fill=orange!8, draw=orange!70, minimum height=0.8cm},
  arr/.style={-{Stealth[length=2.5mm]}, thick, gray!70!black},
  lab/.style={font=\scriptsize, text=gray!50!black}
]
  \node[deci] (d15) {n divisible\\par 15 ?};
  \node[io, right=2.4cm of d15] (r15) {« Thiéboudienne »};
  \node[deci, below=0.9cm of d15] (d3) {n divisible\\par 3 ?};
  \node[io, right=2.4cm of d3] (r3) {« Thié »};
  \node[deci, below=0.9cm of d3] (d5) {n divisible\\par 5 ?};
  \node[io, right=2.4cm of d5] (r5) {« Bou »};
  \node[io, below=0.9cm of d5] (rn) {afficher n};

  \draw[arr] (d15) -- node[lab,above]{oui} (r15);
  \draw[arr] (d15) -- node[lab,left]{non} (d3);
  \draw[arr] (d3) -- node[lab,above]{oui} (r3);
  \draw[arr] (d3) -- node[lab,left]{non} (d5);
  \draw[arr] (d5) -- node[lab,above]{oui} (r5);
  \draw[arr] (d5) -- node[lab,left]{non} (rn);
\end{tikzpicture}

\textit{L'ordre des conditions : on teste d'abord le cas le plus précis (15).}
\end{center}

\begin{lstlisting}[style=js]
function fizzBuzzSenegal() {
  for (let n = 1; n <= 30; n++) {
    if (n % 15 === 0) {
      console.log("Thiéboudienne");
    } else if (n % 3 === 0) {
      console.log("Thié");
    } else if (n % 5 === 0) {
      console.log("Bou");
    } else {
      console.log(n);
    }
  }
}

fizzBuzzSenegal();
\end{lstlisting}

Rappel : \code{n \% 15 === 0} se lit « le reste de la division de \code{n} par 15 est nul », c'est-à-dire « \code{n} est un multiple de 15 ».

\begin{objectifbox}[Exercice 5 --- Palindrome]
Un palindrome est un mot qui se lit pareil à l'endroit et à l'envers (« kayak », « radar », « ressasser »). Écrivez \code{estPalindrome(mot)} qui renvoie un booléen.\\
\emph{Indice :} construisez le mot à l'envers en plaçant chaque lettre \emph{devant} les précédentes, puis comparez avec \code{===}.
\end{objectifbox}

\begin{lstlisting}[style=js]
function estPalindrome(mot) {
  let inverse = "";
  for (const lettre of mot) {
    inverse = lettre + inverse; // on empile à l'envers
  }
  return mot === inverse;
}

console.log(estPalindrome("kayak")); // true
console.log(estPalindrome("Dakar")); // false
\end{lstlisting}

% ===================== AIDE-MÉMOIRE =====================
\section{Aide-mémoire du Lab 5}

\begin{center}
\begin{tabular}{p{5cm} p{8cm}}
\toprule
\textbf{Besoin} & \textbf{Code} \\
\midrule
Brancher le JS & \code{<script src="script.js" defer></script>} \\
Afficher dans la console & \code{console.log(...)} \\
Variable fixe / variable & \code{const x = ...} \quad / \quad \code{let x = ...} \\
Message avec variable & \code{`Texte \$\{variable\}`} \\
Fonction qui renvoie & \code{function f(p) \{ return ...\, \}} \\
Décision & \code{if (...) \{...\} else \{...\}} \quad ou \quad \code{cond ? a : b} \\
Parcourir un tableau & \code{for (const e of tableau) \{...\}} \\
Taille d'un tableau & \code{tableau.length} \\
Reste d'une division & \code{a \% b} \\
\bottomrule
\end{tabular}
\end{center}

% ===================== GIT =====================
\section{Sauvegarder mon travail (routine Git)}

Comme à chaque Lab, on enregistre le travail dans Git et on pose le tag de la séance.

\begin{retenirbox}[Routine Git de fin de Lab --- tag \texttt{v0.5}]
Dans le terminal de VS Code, à la racine du projet :
\begin{lstlisting}[style=html]
git status
git add .
git commit -m "Lab 5 : premières fonctions JavaScript"
git tag v0.5
git push
git push --tags
\end{lstlisting}
Vérifiez ensuite sur GitHub que le commit \emph{et} le tag \code{v0.5} apparaissent bien.
\end{retenirbox}

% ===================== RAPPORT =====================
\begin{attentionbox}[Rapport à produire]
Avant la fin, déposez un rapport court (1 à 2 pages) dans le dossier \code{rapports/} de votre repo, nommé \code{rapport-lab-5-NOM-Prenom.pdf}. Il contient \textbf{quatre} sections :
\begin{enumerate}[leftmargin=*,noitemsep]
  \item \textbf{Ce que j'ai fait} : les fonctions codées (FCFA, moyenne, majorité, FizzBuzz, palindrome).
  \item \textbf{Difficultés rencontrées} : copiez-collez les messages d'erreur qui vous ont bloqué.
  \item \textbf{Comment je les ai résolues} : ce que vous avez cherché, à qui vous avez demandé.
  \item \textbf{Lien vers mon repo} avec le tag : \code{.../jang-afrig/releases/tag/v0.5}.
\end{enumerate}
L'erreur fait partie de l'apprentissage : soyez honnête, c'est le but du rapport.
\end{attentionbox}

% ===================== RÉSUMÉ =====================
\section{Résumé et transition}

Vous venez d'écrire votre \textbf{premier vrai fichier JavaScript} et de coder cinq fonctions qui marchent. Vous savez désormais déclarer des variables, écrire des fonctions, décider avec des conditions et répéter avec \code{for...of}.

Au \textbf{Lab 6}, on monte d'un cran : on apprend à manipuler des \textbf{collections} de données (des tableaux, des objets) avec des outils élégants --- \code{forEach}, \code{map}, \code{filter}. On créera le fichier \code{data.js} contenant nos premiers pays africains : la matière première du futur quiz.

\begin{center}
\textit{Première ampoule allumée. Au Lab 6, on remplit la maison de données.}
\end{center}

\bigskip
\noindent\rule{\textwidth}{0.4pt}
\begin{center}
\small Dr.~El Hadji Bassirou TOURÉ --- UCAD\\
\small Licence 1 DUT1-Informatique --- Fondamentaux du Web --- 2025--2026
\end{center}

\end{document}
